% ============================================================
%  Part IX — Questions
% ============================================================
\gpart{X}{Questions}

% ------------------------------------------------------------
\gsec{W-question words}

\gintro{Question words in German nearly all begin with \textit{w}. Four of them decline for case, in the same pattern as the definite article.}

\begin{gtbl}{colspec={X[0.75,l] X[0.9,l] X[1.35,l]}}
  Word & Meaning & Example \\
  wer     & who \eng{nom.}       & \textbf{Wer} ist das? \\
  wen     & whom \eng{acc.}      & \textbf{Wen} siehst du? \\
  wem     & to whom \eng{dat.}   & \textbf{Wem} gehört das? \\
  wessen  & whose \eng{gen.}     & \textbf{Wessen} Buch ist das? \\
  was     & what                 & \textbf{Was} machst du? \\
  wo      & where \eng{position} & \textbf{Wo} wohnst du? \\
  wohin   & where to             & \textbf{Wohin} gehst du? \\
  woher   & where from           & \textbf{Woher} kommst du? \\
  wann    & when                 & \textbf{Wann} kommst du? \\
  wie     & how                  & \textbf{Wie} geht es dir? \\
  warum   & why                  & \textbf{Warum} lachst du? \\
  wieso   & how come             & \textbf{Wieso} nicht? \\
  weshalb & for what reason      & \textbf{Weshalb} bist du hier? \\
  welcher & which \eng{+ noun}   & \textbf{Welches} Buch? \\
\end{gtbl}

\tcap{wie + another word}

\begin{flattbl}{colspec={X[1,l] X[0.9,l] X[1,l] X[0.9,l]}, vline{3}={1-Z}{0.5pt,solid,Line}, column{3}={font=\small\sffamily\bfseries, fg=Ink}}
  wie viel  & how much  & wie lange & how long \\
  wie viele & how many  & wie oft   & how often \\
  wie weit  & how far   & wie alt   & how old \\
  wie spät  & what time & wie groß  & how big / tall \\
\end{flattbl}

\begin{gnote}
\textit{Wer} declines like \textit{der}: \textit{wer, wen, wem, wessen}. Those
four forms cover every role a person can play in a sentence.
\end{gnote}

% ------------------------------------------------------------
\gsec{Question structure}

\gintro{A yes/no question puts the verb first; a W-question puts the question word first and the verb second. Reported questions send the verb to the end.}

\tcap{Yes/no questions — verb first}

\begin{gtbl}{colspec={X[0.85,l] X[2.15,l]}}
  Type & Example \\
  Statement & Du \textbf{kommst} morgen. \\
  Question  & \textbf{Kommst} du morgen? \\
  With modal & \textbf{Kannst} du morgen kommen? \\
  Perfect    & \textbf{Hast} du das Buch gelesen? \\
\end{gtbl}

\tcap{W-questions — question word first, verb second}

\begin{flattbl}{colspec={X[0.8,l] X[0.6,l] X[1.6,l]}}
  Wann & kommst & du nach Hause? \\
  Warum & hast & du das gemacht? \\
  Mit wem & fährst & du in Urlaub? \\
  Worauf & wartest & du? \\
\end{flattbl}

\tcap{Indirect questions — verb to the end}

\begin{gtbl}{colspec={X[1,l] X[2,l]}}
  Direct & Indirect \\
  Wo wohnt er?      & Ich weiß nicht, wo er \textbf{wohnt}. \\
  Wann kommt sie?   & Sag mir, wann sie \textbf{kommt}. \\
  Kommst du mit?    & Ich frage, \textbf{ob} du \textbf{mitkommst}. \\
\end{gtbl}

\begin{gnote}
An indirect yes/no question has no question word of its own, so German supplies
\textbf{ob} — the equivalent of English \emph{whether}. Never use \textit{wenn}
here.
\end{gnote}

% ------------------------------------------------------------
\gsec{wo(r)- and da(r)- compounds}

\gintro{When a verb is tied to a preposition, German does not say \emph{on what}
or \emph{on it}. It fuses the preposition with \textit{wo-} to ask, and with
\textit{da-} to answer.}

\tcap{Asking: wo(r)- for things, preposition + wen/wem for people}

\begin{gtbl}{colspec={X[0.8,l] X[1.1,l] X[1.1,l]}}
  Verb + prep. & About a thing & About a person \\
  warten auf   & \textbf{Worauf} wartest du? & \textbf{Auf wen} wartest du? \\
  denken an    & \textbf{Woran} denkst du?   & \textbf{An wen} denkst du? \\
  sprechen über & \textbf{Worüber} sprecht ihr? & \textbf{Über wen} sprecht ihr? \\
  sich freuen auf & \textbf{Worauf} freust du dich? & \textbf{Auf wen} freust du dich? \\
\end{gtbl}


\tcap{Answering: da(r)- for things}

\begin{flattbl}{colspec={X[1,l] X[1,l] X[1,l] X[1,l]}, vline{3}={1-Z}{0.5pt,solid,Line}, column{3}={font=\small\sffamily\bfseries, fg=Ink}}
  darauf & on it / for it & daran  & on it / of it \\
  damit  & with it        & davon  & of it, from it \\
  dafür  & for it         & darüber & about it \\
  dazu   & to it          & darunter & under it \\
\end{flattbl}

\begin{gnote}
The \textit{r} appears only when the preposition begins with a vowel:
\textit{wo\textbf{r}auf}, \textit{da\textbf{r}an}, \textit{wo\textbf{r}über},
against \textit{womit}, \textit{davon}, \textit{wofür}. People never take these
compounds — for a person the preposition keeps its pronoun:
\textit{auf \textbf{ihn}}, \textit{mit \textbf{ihr}}.
\end{gnote}
